% Deducción propia a partir de la variación ya establecida. No es un teorema del paper LL.
\section{Estructura holográfica generalizada}
\SetRoadmapStage{7}
\setcounter{holografiaSavedEquation}{\value{equation}}
\begingroup
\setcounter{equation}{0}
\renewcommand{\theequation}{G\arabic{equation}}
\renewcommand{\theHequation}{holografia-gen.\arabic{equation}}
\MapaHolografiaGen{0}
\begin{frame}[label=holo-gen-01]{Hipótesis y deducción del defecto de escalamiento}
\hypertarget{holo-gen-04}{}
\footnotesize\CompactDisplaySpacing
Partimos de $L(g,R,\phi,\nabla\phi)$ sin derivadas de $R$, con conexión de Levi--Civita. Suponemos que $\mathscr L=\sqrt{-g}L$ tiene peso $\kappa$ bajo $g_{ab}\mapsto\lambda g_{ab}$ constante, con $\phi$ y sus derivadas coordenadas fijas. Los pesos distintos se tratan por sectores.
\[
\delta\mathscr L=\sqrt{-g}(E_{ab}\delta g^{ab}+E_\phi\delta\phi)
+\partial_j(\sqrt{-g}\Theta^j),
\]
\begin{equation}
\Theta^j=2P^{ibjd}\nabla_b\delta g_{di}-2\delta g_{di}\nabla_cP^{ijcd}+\Pi^j\delta\phi.
\label{eq:hg-boundary}
\end{equation}
Elegimos $\delta g_{ab}=\epsilon g_{ab}$ y $\delta\phi=0$: entonces $\delta g^{ab}=-\epsilon g^{ab}$, $\nabla_b\delta g_{di}=0$ y $\Theta^j=-2\epsilon g_{di}\nabla_cP^{ijcd}$. Sustituir y dividir por $\epsilon$ da
\[
\kappa\mathscr L=-\sqrt{-g}g^{ab}E_{ab}-2\partial_j(\sqrt{-g}g_{di}\nabla_cP^{ijcd}).
\]
Definimos la divergencia residual como
\begin{equation}
\boxed{\Delta_{\rm H}\equiv\kappa\mathscr L+\sqrt{-g}g^{ab}E_{ab}
=-2\partial_j(\sqrt{-g}g_{di}\nabla_cP^{ijcd}).}
\label{eq:hg-defect}
\end{equation}
$\nabla P=0$ implica $\Delta_{\rm H}=0$, no a la inversa. Fijar $\delta\phi=0$ no elimina el término métrico con $\nabla P$. Es una deducción propia fuera de las soluciones, no un teorema atribuido al artículo LL.
\end{frame}
\MapaHolografiaGen{1}
\begin{frame}[label=holo-gen-02]{Separación exacta y orden de las derivadas}
\footnotesize
Si un sector $L_R$ es homogéneo de grado $m\ge1$ en curvatura, Euler da
$P^{abcd}R_{abcd}=mL_R$. Con $Q=P/m$, \eqref{eq:holo-split-general} implica
\begin{equation}
\mathscr L_R=\mathscr L_{\rm vol}+\mathscr L_{\rm sup},\qquad
\mathscr L_{\rm vol}=\mathscr L_{\rm vol}^{(0)}+\mathscr L_{\nabla P},
\label{eq:hg-split}
\end{equation}
\[
\begin{aligned}
\mathscr L_{\rm sup}&=\frac2m\partial_c(\sqrt{-g}P_a{}^{bcd}\Gamma^a_{bd}),\\
\mathscr L_{\rm vol}^{(0)}&=\frac2m\sqrt{-g}P_a{}^{bcd}\Gamma^a_{de}\Gamma^e_{bc},\\
\mathscr L_{\nabla P}&=-\frac2m\sqrt{-g}\Gamma^a_{bd}\nabla_cP_a{}^{bcd}.
\end{aligned}
\]
\textbf{El volumen incluye la corrección.} Usar solo $\mathscr L_{\rm vol}^{(0)}$ fuera de LL no reconstruye la acción.
\medskip

Si $P$ depende de $R$, $\nabla P$ puede contener $\partial^3g$.
La fórmula de segundo gradiente no basta entonces.
Para $Q_\beta$, lineal en curvatura, $P_\beta=P_\beta(g,u)$ y el volumen contiene hasta $\partial g$ y $\partial^2\phi$, sin $\partial^2g$.

Los términos sin curvatura pueden asignarse al volumen. La descomposición y sus partes dependen de las coordenadas.
\end{frame}
\MapaHolografiaGen{2}
\begin{frame}[label=holo-gen-03]{Operador de escalamiento del volumen}
\footnotesize
Con $q^A=g_{ab}$ y manteniendo todos los argumentos escalares fijos, para un volumen de segundo gradiente definimos
\[
M_A^{ij}=\frac{\partial\mathscr L_{\rm vol}}{\partial q^A_{,ij}},\qquad
p_A^i=\frac{\partial\mathscr L_{\rm vol}}{\partial q^A_{,i}}-\partial_jM_A^{ij}.
\]
Las dos integraciones por partes de \eqref{eq:holo-second-variation}, con $\delta q^A=\epsilon q^A$, dan
\begin{equation}
\mathcal H_g[\mathscr L_{\rm vol}]
\equiv\partial_i(q^Ap_A^i+q^A_{,j}M_A^{ij})
=\kappa\mathscr L_{\rm vol}-g_{ab}\frac{\delta\mathscr L_{\rm vol}}{\delta g_{ab}}.
\label{eq:hg-operator}
\end{equation}
Si aparecen terceros gradientes, sea $T_A^{ijk}=\partial\mathscr L_{\rm vol}/\partial q^A_{,ijk}$, simétrico. Una integración adicional produce
\[
M_A^{ij}=\frac{\partial\mathscr L_{\rm vol}}{\partial q^A_{,ij}}-\partial_kT_A^{ijk},\quad
p_A^i=\frac{\partial\mathscr L_{\rm vol}}{\partial q^A_{,i}}-\partial_jM_A^{ij},
\]
\[
\mathcal H_g=\partial_i(q^Ap_A^i+q^A_{,j}M_A^{ij}+q^A_{,jk}T_A^{ijk}).
\]
La igualdad de escalamiento de \eqref{eq:hg-operator} se conserva. Para $Q_\beta$, $M=T=0$ en la variación métrica, aunque haya segundos gradientes escalares.
\end{frame}
\MapaHolografiaGen{3}

\begin{frame}[label=holo-gen-05]{Identidad volumen--superficie con defecto}
Para una separación homogénea $\mathscr L=\mathscr L_{\rm vol}+\mathscr L_{\rm sup}$, la superficie es una divergencia y
\[
\frac{\delta\mathscr L_{\rm vol}}{\delta g_{ab}}
=\frac{\delta\mathscr L}{\delta g_{ab}}=-\sqrt{-g}E^{ab}.
\]
Por \eqref{eq:hg-operator},
\begin{align*}
\mathcal H_g[\mathscr L_{\rm vol}]
&=\kappa\mathscr L_{\rm vol}+\sqrt{-g}g^{ab}E_{ab}\\
&=\kappa\mathscr L_{\rm vol}+\Delta_{\rm H}-\kappa\mathscr L\\
&=-\kappa\mathscr L_{\rm sup}+\Delta_{\rm H}.
\end{align*}
Por tanto,
\begin{equation}
\boxed{\kappa\mathscr L_{\rm sup}=-\mathcal H_g[\mathscr L_{\rm vol}]+\Delta_{\rm H}.}
\label{eq:hg-identity}
\end{equation}
Para un término LL geométrico, $\kappa=D/2-m$ y $\Delta_{\rm H}=0$: recuperamos \eqref{eq:holo-LL-identity}.

Si $\kappa=0$, queda $\mathcal H_g=\Delta_{\rm H}$ y la identidad no reconstruye la superficie por división.
\end{frame}
\MapaHolografiaGen{4}
\begin{frame}[label=holo-gen-06]{Superficie efectiva: significado y límite}
Para $\kappa\ne0$ podemos reorganizar \eqref{eq:hg-identity} mediante
\begin{equation}
\widetilde{\mathscr L}_{\rm sup}
=\mathscr L_{\rm sup}-\frac{\Delta_{\rm H}}{\kappa}
=\mathscr L_{\rm sup}+\frac2\kappa\partial_j(\sqrt{-g}g_{di}\nabla_cP^{ijcd}),
\label{eq:hg-effective}
\end{equation}
\[
\kappa\widetilde{\mathscr L}_{\rm sup}=-\mathcal H_g[\mathscr L_{\rm vol}].
\]
\textbf{Aquí el operador actúa sobre el volumen original.} La pareja
$\mathscr L_{\rm vol}+\widetilde{\mathscr L}_{\rm sup}$ no suma la acción original, sino $\mathscr L-\Delta_{\rm H}/\kappa$.
Para conservar la misma densidad debemos definir también
\[
\widetilde{\mathscr L}_{\rm vol}=\mathscr L_{\rm vol}+\Delta_{\rm H}/\kappa.
\]
Como $\Delta_{\rm H}$ es una divergencia homogénea, su derivada de Euler es cero y $\mathcal H_g[\Delta_{\rm H}/\kappa]=\Delta_{\rm H}$. Luego
\[
\kappa\widetilde{\mathscr L}_{\rm sup}
=-\mathcal H_g[\widetilde{\mathscr L}_{\rm vol}]+\Delta_{\rm H}.
\]
La redefinición no elimina el defecto de una separación consistente. Es una reescritura útil, no una prueba de que la teoría sea LL.
\end{frame}
\MapaHolografiaGen{5}

\begin{frame}[label=holo-gen-07]{Sector $Q_\beta$: peso y contracción de $P_\beta$}
\footnotesize\CompactDisplaySpacing
Para $Q_\beta=\ell^2\beta_0(3R_{ab}u^au^b-RX)$, $u_a=\partial_a\phi$ y $X=u_au^a$, el escalamiento a $\phi$ fijo da
\[
u_a\mapsto u_a,\quad u^a\mapsto\lambda^{-1}u^a,\quad
R_{ab}\mapsto R_{ab},\quad R,X\mapsto\lambda^{-1}(R,X).
\]
Así $Q_\beta\mapsto\lambda^{-2}Q_\beta$, aunque es lineal en curvatura:
\begin{equation}
\kappa_\beta=D/2-2\qquad (m_\beta=1).
\label{eq:hg-beta-weight}
\end{equation}
La contracción del momento ya calculado se obtiene término a término:
\[
\begin{aligned}
g_{bc}P_\beta^{abcd}
=\frac{\ell^2\beta_0}{4}\big[&3(u^au^d+u^au^d-Du^au^d-Xg^{ad})\\
&-2X(g^{ad}-Dg^{ad})\big].
\end{aligned}
\]
Agrupando,
\begin{equation}
g_{bc}P_\beta^{abcd}
=\frac{\ell^2\beta_0}{4}[-3(D-2)u^au^d+(2D-5)Xg^{ad}].
\label{eq:hg-beta-contraction}
\end{equation}
El peso métrico no se obtiene contando solo curvaturas cuando hay gradientes escalares.
\end{frame}
\begin{frame}[label=holo-gen-08]{Sector $Q_\beta$: divergencia y superficie efectiva}
\footnotesize\CompactDisplaySpacing
Las dos antisimetrías de $P$ permiten escribir
$g_{di}\nabla_cP^{ijcd}=g_{bc}\nabla_dP^{jbcd}$.
Por compatibilidad métrica, contraer y derivar conmutan. En $D=3$,
\[
C_\beta^a\equiv g_{bc}\nabla_dP_\beta^{abcd}
=\frac{\ell^2\beta_0}{4}[-3\nabla_d(u^au^d)+\nabla^aX].
\]
Al introducirlo en \eqref{eq:hg-defect},
\begin{equation}
\Delta_{{\rm H},\beta}
=\frac{\ell^2\beta_0}{2}\partial_a\!\left\{\sqrt{-g}
[3\nabla_d(u^au^d)-\nabla^aX]\right\}.
\label{eq:hg-beta-defect}
\end{equation}
Como $\kappa_\beta=-1/2$, \eqref{eq:hg-effective} da
\begin{equation}
\widetilde{\mathscr L}_{{\rm sup},\beta}
=\mathscr L_{{\rm sup},\beta}+\ell^2\beta_0\partial_a\!\left\{\sqrt{-g}
[3\nabla_d(u^au^d)-\nabla^aX]\right\}.
\label{eq:hg-beta-effective}
\end{equation}
\[
-\tfrac12\widetilde{\mathscr L}_{{\rm sup},\beta}
=-\mathcal H_g[\mathscr L_{{\rm vol},\beta}].
\]
La fórmula conserva el volumen original, incluida $\mathscr L_{\nabla P}$. En $D=4$, $\kappa_\beta=0$: esta división no está permitida.
\end{frame}
\begin{frame}[label=holo-gen-09]{Chequeo sobre el ansatz estático}
\footnotesize\CompactDisplaySpacing
Como comprobación, usamos $\dd s^2=-f(r)\dd t^2+\dd r^2/f(r)+r^2\dd\theta^2$ y $u_a=(0,0,p)$.
Entonces $\sqrt{-g}=r$, $X=p^2/r^2$ y $u^\theta=p/r^2$.
\[
\nabla_du^d=0,\quad \Gamma^r_{\theta\theta}=-rf,\quad
\nabla_d(u^ru^d)=\Gamma^r_{\theta\theta}(u^\theta)^2=-\frac{fp^2}{r^3},\quad
\nabla^rX=-\frac{2fp^2}{r^3}.
\]
Los otros componentes de $C_\beta^a$ se anulan. Por tanto,
\[
C_\beta^r=\frac{\ell^2\beta_0fp^2}{4r^3},\qquad
\Delta_{{\rm H},\beta}=-2\partial_r(rC_\beta^r)
=-\frac{\ell^2\beta_0p^2}{2}\frac{\dd}{\dd r}\left(\frac f{r^2}\right).
\]
La corrección superficial es
\[
\widetilde{\mathscr L}_{{\rm sup},\beta}-\mathscr L_{{\rm sup},\beta}
=-\ell^2\beta_0p^2\left(\frac{f'}{r^2}-\frac{2f}{r^3}\right).
\]
Si $f\propto r^2$, el defecto se anula en este ansatz. Sin embargo, la componente calculada en la discusión, $\nabla_aP_\beta^{attr}=\ell^2\beta_0p^2/(4r^3)$, puede ser no nula.
Es un chequeo sobre una geometría, no una identidad LL ni una imposición de las ecuaciones de campo.
\end{frame}
\begin{frame}[label=holo-gen-10]{Resumen: identidad con defecto y frontera}
\hypertarget{nav-map-7-6}{}
\footnotesize
Para cualquier separación $\mathscr L=\mathscr L_{\rm vol}+\partial_aV^a$,
\[
\delta\mathscr L_{\rm vol}=\sqrt{-g}(E_{ab}\delta g^{ab}+E_\phi\delta\phi)
+\partial_a\vartheta_{\rm vol}^a
\]
implica, con una elección compatible de corrientes,
\[
\sqrt{-g}\Theta^a=\vartheta_{\rm vol}^a+\delta V^a
\quad\text{hasta una corriente de divergencia idénticamente nula}.
\]
Por eso $\delta V^a$ no tiene que coincidir por sí sola con $\sqrt{-g}\Theta^a$.
Si $V^a$ cambia por $2\sqrt{-g}C^a/\kappa$, con $C^a=g_{bc}\nabla_dP^{abcd}$, la frontera de la acción cambia por la variación de esa corrección, salvo compensación en el volumen.
\begin{block}{Qué queda establecido}
La homogeneidad produce una identidad con defecto $\Delta_{\rm H}$. El caso $Q_\beta$ tiene peso $D/2-2$ y la corrección explícita anterior.
La superficie efectiva reorganiza esa identidad, pero no prueba una nueva dualidad ni un principio de borde de tipo LL.
\end{block}
Si $\delta\phi\ne0$ debe conservarse $\Pi_\beta^a\delta\phi$. Las condiciones de contorno y los posibles contraterminos requieren un análisis variacional propio.
\global\mapcompletedfooter=6\relax
\end{frame}
\setcounter{equation}{\value{holografiaSavedEquation}}
\endgroup
